\documentclass[a4paper, 12pt,table, hidelinks, DIV=calc]{extarticle} % что то меняется , не понятно article or extarticle...
%\usepackage[T2A]{fontenc}
\usepackage{fontspec}
\usepackage[english,russian]{babel}
\usepackage{graphicx}

\setmainfont{Arial Narrow}
\usepackage{unicode-math} % Задать шрифт для математических формул
\setmathfont{Stix Two Math}
%\setmathfont{ArialNarrow}
%\setmathfont{Cambria Math}

% для KOMA script чтобы в документе менять формат А3 А4 и ориентацию
\usepackage{typearea}

\usepackage[
a4paper,
includehead,
includefoot,
nomarginpar,
top=6mm,
bottom=6mm,
left=22mm,
right=10mm,
headheight=14.5pt,
headsep=3mm,        % Добавить headsep
footskip=7mm,
%textheight=276mm    % Должен удовлетворять формуле
]{geometry}

%\savegeometry{standard}
%\setlength{\topmargin}{-27.5mm}
%\setlength{\headsep}{3mm} % Устанавливаем отступ между колонтитулом и текстом 10pt !!! Сдвигает а4 landscape
%\setlength{\topskip}{10pt} % Установка вертикального промежутка между верхним краем текста и колонтитулом на 0pt
%\setlength{\textheight}{276mm}

%\usepackage{fancyhdr}

%\usepackage{lmodern}

%\usepackage{longtable}
%\usepackage{amsfonts}

\usepackage{xcolor} % ПРИ КОМПИЛЯЦИИ В ЛИНУКС УБРАТЬ ЭТУ СТРОКУ, да вообще ее убрать? или нет???
\usepackage{colortbl}

%\usepackage{array}

\usepackage{hyperref}

\graphicspath{{./images/}}

\usepackage{fontspec}
\setsansfont[
        Ligatures=TeX,
        BoldFont={Arial Narrow Bold},
        ItalicFont={Arial Narrow Italic},
        BoldItalicFont={Arial Narrow Bold Italic}
    ]{Arial Narrow}
\renewcommand\familydefault{\sfdefault}
%\setmainfont{Arial Narrow}
\usepackage{fancyhdr}

\usepackage{parskip}% для абзацев и отступа
\setlength{\parindent}{1cm}
\setlength{\parskip}{1pt}
\setlength{\intextsep}{0mm} % отступ между плавающими элементами


\usepackage{indentfirst}

% Определяем цвета фирмы старые
%\definecolor{unired}{RGB}{247, 148, 29} 
%\definecolor{uniblue}{RGB}{30, 73, 115}
%\definecolor{unigray}{RGB}{88, 89, 91}

% Определяем новые цвета фирмы (unitel guide 23.10.24)
\definecolor{unired}{RGB}{239, 134, 55} 
\definecolor{uniblue}{RGB}{21, 73, 140}
\definecolor{unigray}{RGB}{144, 143, 143}
\definecolor{uniblack}{RGB}{20, 20, 20}

% Определяем новые цвета фирмы (unitel guide 30.03.26)
\definecolor{uniwhite}{RGB}{255, 255, 255} 
\definecolor{uniblack}{RGB}{20, 20, 20}
\definecolor{unigreen}{RGB}{0, 203, 98}
\definecolor{unidarkgreen}{RGB}{6, 150, 75}
\definecolor{unigray}{RGB}{100, 100, 100}

\usepackage{longtable,booktabs}
\usepackage{calc} % for calculating minipage widths

% Это для таблицы, чтобы жестко ширину столбцов выставлять НЕ ТРЕБУЕТСЯ???
%\usepackage{array}
%\newcolumntype{L}[1]{>{\raggedright\let\newline\\\arraybackslash\hspace{0pt}}m{#1}}
%\newcolumntype{C}[1]{>{\centering\let\newline\\\arraybackslash\hspace{0pt}}m{#1}}
%\newcolumntype{R}[1]{>{\raggedleft\let\newline\\\arraybackslash\hspace{0pt}}m{#1}}

%\usepackage{bigstrut}

\usepackage[labelsep=endash]{caption} % выравнить надписи к рисункам тире вместо двоеточия
\DeclareCaptionLabelFormat{myformat}{\textls[75]{#1 #2}} % Здесь устанавливается разрежение текста для надписей Рисунок и Таблица
\captionsetup[figure]{name=Рисунок, justification=centering, labelformat=myformat} 
\captionsetup[table]{justification=justified, labelformat=myformat} % Убираем отступ при переносе строки в наименовании таблицы
\setlength{\belowcaptionskip}{0mm} % отступ после подписи
%\setlength{\abovecaptionskip}{0mm} % отступ до подписи

% нумерация таблиц и рисунков ИСПОЛЬЗУЕТСЯ!!!
% * код для нумерации вида 1.1.1
%\counterwithin*{figure}{subsection}
%\counterwithin*{table}{subsection}
%\renewcommand{\thefigure}{%
  %\ifnum\value{subsection}>0
    %\thesubsection.%
  %\else
    %\thesection.%
  %\fi
  %\arabic{figure}%
%}
%\renewcommand{\thetable}{%
  %\ifnum\value{subsection}>0
    %\thesubsection.%
  %\else
    %\thesection.%
  %\fi
  %\arabic{table}%
%}
% * код для нумерации вида 1.1
%\counterwithin*{figure}{section}
%\counterwithin*{table}{section}
%\renewcommand{\thefigure}{%
  %\ifnum\value{subsection}>0
    %\thesection.% 
  %\else
    %\thesection.% 
  %\fi
  %\arabic{figure}%
%}
%\renewcommand{\thetable}{%
  %\ifnum\value{subsection}>0
    %\thesection.% 
  %\else
    %\thesection.%
  %\fi
  %\arabic{table}%  
%}
% Оптимизированная нумерация вида 1.1 с нормальным закрытием скобок ниже
\counterwithin*{figure}{section}
\counterwithin*{table}{section}
\renewcommand{\thefigure}{%
  \thesection.%
  \arabic{figure}%
}
\renewcommand{\thetable}{%
  \thesection.%
  \arabic{table}%
}

%%%%%%%%%%%%%%% ЭТОТ КОД ДЕЛАЕТ ТОЖЕ САМОЕ? %%%%%%%%%%%%%%%%%%%
%\usepackage{chngcntr}
%\counterwithin{figure}{subsection}
%\counterwithin{table}{subsection}
%\renewcommand{\thefigure}{\thesubsection.\arabic{figure}}
%\renewcommand{\thetable}{\thesubsection.\arabic{table}}
%%%%%%%%%%%%%%%%%%%%%%%%%%%%%%%%%%%%%%%%%%%%%%%%%%%%%%%%%%%%%%%

\usepackage{tabularx}
%\usepackage[table]{xcolor}
\usepackage{multirow} % Для центрирования содержимого ячейки по вертикали - ИСПОЛЬЗУЕТСЯ!!!
\usepackage{amsmath} % для вывода формул
\usepackage{eqexpl} % для вывода пояснений к формулам
\eqexplSetDelim{--}
\eqexplSetIntro{где}

\usepackage{titlesec}
\usepackage{appendix}
%\titlespacing{\chapter}{0pt}{-30pt}{12pt} как пример
\titlespacing{\section}{30pt}{0pt}{7pt}
\titlespacing{\subsection}{30pt}{6pt}{5pt}
\titlespacing{\subsubsection}{30pt}{3pt}{3pt}

%\usepackage[bitstream-charter]{mathdesign} % проверить когда формулы будут в тексте
\uchyph=0 % отключаем перенос в словах с заглавной буквы - чтобы ГЗоткл не переносил например

% Заменяем точки на тире в перечислении (списке)
% - дефис, -- тире, --- длинное тире
\usepackage{enumitem}

\setlist[itemize]{label=---}
%\setlist[itemize]{label=$\circ$} % окружность
%%%%%%%%%%%%%%%%%%%%%%%%

\usepackage{float}

% не допускает в конце странцы заголовок и одна или две строки , переносит на след страницу
% необходимо указывать непосредственно в разделе
\usepackage{needspace} %ИСПОЛЬЗУЕТСЯ!!!
%%%%%%%%%%%%%%%%%%%%%%%%

\usepackage{tabularray}
\usepackage{tabularx}

\usepackage{tocloft}  % точки для глав в оглавлении ИСПОЛЬЗУЕТСЯ!!!
%\renewcommand{\cftsecfont}{\normalfont\scshape} 
\usepackage{tocvsec2} % для подавления отображения подразделов в toc для приложений, не знаю, если подразделы появятся
\renewcommand{\cftsecleader}{\cftdotfill{\cftdotsep}} % точки для разделов в оглавлении ИСПОЛЬЗУЕТСЯ!!!
\renewcommand{\cftsecpagefont}{\normalfont} % задаем не жирный формат номеров страниц в оглавлении, если строчку убрать - будет жирный

\usepackage{pdflscape}
%\titleformat{\section}{\normalfont\Large\bfseries}{\thesection}{1em}{} % 01.11.24 отключил
%\titleformat{\section}{\normalfont\large\bfseries\MakeUppercase}{\thesection}{1em}{} % Заголовки разделов большими буквами

%\usepackage{lscape} % УДАЛИТЬ???
%\usepackage{lipsum}
\usepackage[absolute]{textpos}
%\usepackage{etoolbox}

\usepackage{pdfpages} % попробуем вставить векторный рисунок

%\titlespacing{\subsection}{0pt}{\parskip}{-\parskip} % отступ между сабсекцией и лонгтейбл?
%\titlespacing{command}{left spacing}{before spacing}{after spacing}[right]

%\titlespacing\section{0pt}{12pt plus 4pt minus 2pt}{0pt plus 2pt minus 2pt}
%\titlespacing\subsection{0pt}{12pt plus 4pt minus 2pt}{0pt plus 2pt minus 2pt}
%\titlespacing\subsubsection{0pt}{12pt plus 4pt minus 2pt}{0pt plus 2pt minus 2pt}

%\titlespacing{\section}{0pt}{\parskip}{-\parskip}
%\titlespacing{\subsection}{0pt}{\parskip}{-\parskip}
%\titlespacing{\subsubsection}{0pt}{\parskip}{-\parskip}

\usepackage{fontawesome} % Подключить пакет шрифтов fontawesome

\usepackage{microtype} % для разрежения текста

\usepackage{etoolbox} % для проверки логических условий

% ВОТ ЭТА ХРЕНЬ НИЖЕ ВМЕСТЕ С ПАКЕТАМИ etoolbox и tocloft делает заголовки Разделов большими буквами
% https://latex.org/forum/viewtopic.php?t=24404
\makeatletter
\patchcmd{\l@section}
  {\cftsecfont #1}%   search pattern
  {\cftsecfont {#1}}% replace by
  {}%                  success
  {}%                  failure
\makeatother
\renewcommand\cftsecfont{\MakeUppercase}
%%%%%%%%%%%%%%%%%%%%%%%%%%%%%%%%%%%%%%%%%%%%%%%%%

% Делаем заголовки подразделов большими буквами ВЫГЛЯДИТ НЕ ОЧЕНЬ, ГОСТ 2.105-2019 ,п.6.2.5
%\makeatletter
%\patchcmd{\l@subsection}
  %{\cftsubsecfont #1}%   search pattern
  %{\cftsubsecfont {#1}}% replace by
  %{}%                  success
  %{}%                  failure
%\makeatother
%\renewcommand\cftsubsecfont{\MakeUppercase}

%%%%%%%%% В ЭТОМ БЛОКЕ ПОДКЛЮЧЕНЫ БЛОКИ ДЛЯ ПРИЛОЖЕНИЯ КАРТА ЗАКАЗА %%%%%%%%%%%
\usepackage{tikz}% для рисования если нет рисования то удалить
\usepackage{makecell}
%%%%%%%%%%%%%%%%%%%%%%%%%%% КОНЕЦ БЛОКА %%%%%%%%%%%%%%%%%%%%%%%%%%%%%%%%%%%%%%%

%%%%%%%% В ЭТОМ БЛОКЕ ДЕЛАЕМ ЧТОБЫ СНОСКА БЫЛА СО СКОБКОЙ, НАПРИМЕР 1) 2) БЕЗ НЕГО СНОСКА ОБЫЧНОЕ ЧИСЛО %%%%%
\usepackage[perpage]{footmisc}
\renewcommand{\thefootnote}{\arabic{footnote})}
% Добавляем отступ между текстом и сносками
\setlength{\footnotesep}{10pt} % измените 12pt на нужное вам значение
%%%%%%%%%%%%%%%%%%%%%%%%%%% КОНЕЦ БЛОКА %%%%%%%%%%%%%%%%%%%%%%%%%%%%%%%%%%%%%%%

% Устанавливаем толщину линий
%\setlength{\arrayrulewidth}{0.5pt} % Толщина линий таблицы

\usepackage{hhline} % чтобы рисовать по одной линии в таблицах, соответственно тэг \hhline{} должен применяться

%\setcounter{secnumdepth}{3} % Нумерация до разделов, подразделов и подподразделов
%\setcounter{tocdepth}{2}     % Запись в содержание только разделов и подразделов

\usepackage{ulem} % 29.11.24 - чтобы подчеркивание в ссылках более менее нужно
\normalem % 29.11.24 - чтобы от верхнего пакета не ломалось \emph

\usepackage{placeins} % чтобы ограничивать float рисунков с тегом [!h]

\usepackage{eso-pic}